% This file compiles standalone.  The guard below also permits a future lite
% wrapper to define \litemode and provide the shared lite setup.
\ifx\litemode\undefined\documentclass{article}\fi
\usepackage[margin=3cm]{geometry}
\usepackage{amsmath,amssymb}
\usepackage{mathtools}
\usepackage{amsthm}
\usepackage{enumitem}
\usepackage{hyperref}
\usepackage{cleveref}
\crefname{property}{property}{properties}
\Crefname{property}{Property}{Properties}
\crefname{assumption}{assumption}{assumptions}
\Crefname{assumption}{Assumption}{Assumptions}

\pagestyle{plain}

% Theorem environments
\theoremstyle{plain}
\newtheorem{theorem}{Theorem}
\newtheorem*{theorem*}{Theorem}
\newtheorem{corollary}{Corollary}
\newtheorem{lemma}{Lemma}
\newtheorem{proposition}{Proposition}
\newtheorem{property}[theorem]{Property}

\theoremstyle{definition}
\newtheorem{definition}{Definition}

\theoremstyle{remark}
\newtheorem{remark}[theorem]{Remark}
\newtheorem{assumption}[theorem]{Assumption}

\IfFileExists{lite_full_setup.tex}{% Shared lite/full setup for protocol-spec .tex files.
%
% Files under full/ that include this snippet compile in two modes:
%   - Full (default): all content rendered.
%   - Lite: when \litemode is defined before \documentclass (typically by a
%     wrapper .tex file at the repo root), lemmas, corollaries, propositions,
%     properties, proofs, remarks, assumptions, and conditions are dropped
%     via the comment package. Algorithms, theorem statements, definitions,
%     and surrounding prose remain.
%
% Place \input of this file in each spec's preamble AFTER all \newtheorem
% declarations (or just before \begin{document}). The exclusions use
% \@ifundefined so it's safe even when a file defines only a subset.
%
% Path resolution: \IfFileExists tries the local copy first (when cwd is
% full/), then falls back to full/lite_full_setup.tex (when cwd is the repo
% root, e.g. while compiling a top-level _lite.tex wrapper).
\newif\iffull
\ifx\litemode\undefined\fulltrue\else\fullfalse\fi
\iffull\else
  \usepackage{comment}
  \makeatletter
  \@ifundefined{lemma}{}{\excludecomment{lemma}}
  \@ifundefined{corollary}{}{\excludecomment{corollary}}
  \@ifundefined{proposition}{}{\excludecomment{proposition}}
  \@ifundefined{property}{}{\excludecomment{property}}
  \@ifundefined{proof}{}{\excludecomment{proof}}
  \@ifundefined{remark}{}{\excludecomment{remark}}
  \@ifundefined{assumption}{}{\excludecomment{assumption}}
  \@ifundefined{condition}{}{\excludecomment{condition}}
  \makeatother
\fi
}{%
  \IfFileExists{full/lite_full_setup.tex}{% Shared lite/full setup for protocol-spec .tex files.
%
% Files under full/ that include this snippet compile in two modes:
%   - Full (default): all content rendered.
%   - Lite: when \litemode is defined before \documentclass (typically by a
%     wrapper .tex file at the repo root), lemmas, corollaries, propositions,
%     properties, proofs, remarks, assumptions, and conditions are dropped
%     via the comment package. Algorithms, theorem statements, definitions,
%     and surrounding prose remain.
%
% Place \input of this file in each spec's preamble AFTER all \newtheorem
% declarations (or just before \begin{document}). The exclusions use
% \@ifundefined so it's safe even when a file defines only a subset.
%
% Path resolution: \IfFileExists tries the local copy first (when cwd is
% full/), then falls back to full/lite_full_setup.tex (when cwd is the repo
% root, e.g. while compiling a top-level _lite.tex wrapper).
\newif\iffull
\ifx\litemode\undefined\fulltrue\else\fullfalse\fi
\iffull\else
  \usepackage{comment}
  \makeatletter
  \@ifundefined{lemma}{}{\excludecomment{lemma}}
  \@ifundefined{corollary}{}{\excludecomment{corollary}}
  \@ifundefined{proposition}{}{\excludecomment{proposition}}
  \@ifundefined{property}{}{\excludecomment{property}}
  \@ifundefined{proof}{}{\excludecomment{proof}}
  \@ifundefined{remark}{}{\excludecomment{remark}}
  \@ifundefined{assumption}{}{\excludecomment{assumption}}
  \@ifundefined{condition}{}{\excludecomment{condition}}
  \makeatother
\fi
}{}%
}

\title{Fresh Simplex with Height Filter and Timeouts}
\author{}

\begin{document}
\maketitle

\begin{abstract}
We describe Fresh Simplex and prove accountable safety and recovery after a
period of asynchrony.  Quorum checks at a given height use fixed integer
weights.  A vote for the current-height target and a timeout both count toward
height advancement, while an empty vote does not.  The chain state records the
target and therefore needs only one target-participation bit and one timeout
bit per validator.  Signed votes are still kept as evidence for E1 and E2.

Let~\(H\) be nonjustifiable on every valid branch, including branches revealed
after recovery begins.  Once the protocol advances from~\(H\) through
\(H+1\) to~\(H+2\), Goldfish confirmation and the stabilization gadget keep
the confirmed available chain stable.  To exclude an older justification at
height~\(H+2\) or above, \(H+2\) must be greater than every height previously
reached or signed at by an honest validator.

Available-chain stability is not by itself a finality result.  Finality also
requires a later ordinary height at which no honest validator has previously
voted.  The votes at that height must be included on time, the resulting
justification must be known before the following finality votes, and it must
remain the latest justification until those votes are included.  Under these
conditions the ordinary height is justified and then finalized.
\end{abstract}

\section{Protocol overview}\label{sec:overview}

The finality state machine and the fork-choice rule are defined separately.
Goldfish gives the available-chain confirmation.  The stabilization gadget
uses latest head votes to keep later fork-choice walks on that chain, and the
finality gadget advances height and finalizes justified blocks.  The
\(G_0\) test defines a safe-confirmed prefix for target and timeout votes; it
does not alter Goldfish confirmation.  The recovery proof uses compatible
Goldfish confirmations during the recovery period, but does not require
honest validators to have identical vote sets or denominators.

\section{Model and quorum weights}\label{sec:model}

\begin{definition}[Validator weights]\label{def:validator-weights}
At every finality height \(h\), let \(V_h\) be the validator set and let
\(w_h(i)\) be validator~\(i\)'s positive integer weight.  Write
\[
  W_h=\sum_{i\in V_h}w_h(i),
  \qquad
  q_h=\left\lceil\frac{2W_h}{3}\right\rceil,
  \qquad
  \iota_h=2q_h-W_h.
\]
For \(S\subseteq V_h\), write \(w_h(S)=\sum_{i\in S}w_h(i)\).  A quorum at
height~\(h\) is a set \(Q\subseteq V_h\) with \(w_h(Q)\geq q_h\).
All threshold checks use integer weights or the equivalent cross-multiplication
\(3w_h(Q)\geq 2W_h\); no rounded \emph{adjusted denominator} is part of the
protocol.
\end{definition}

\begin{assumption}[Validator weights at each height]\label{ass:height-weights}
For every finality height~\(h\), all valid chains use the same \(V_h\),
\(w_h\), \(W_h\), and \(q_h\).  Target, timeout, and finality votes at
height~\(h\) are all counted using these weights.  Stabilization-gadget head
messages use the validator weights for their round, as specified in
Assumption~\ref{ass:recovery-weights}.  A proposer pointer contains only a
block root.  Its previous-round support is evaluated using the weights for
that round, regardless of the receiver's current finality height.  The
Byzantine weight \(b_h\) satisfies
\[
  b_h < \iota_h.
\]
The validator set and weights for height~\(h\) are fixed once that height is
reachable and are the same on all valid chains in the weak-subjectivity
domain.
\end{assumption}

This assumption holds for a fixed validator set.  It also holds when validator
changes for a new height are derived from previously finalized state.  It does
not hold if competing block states can choose different weights for the same
height.  Quorum intersection is used only for quorums counted with the same
weights.

\begin{lemma}[Weighted quorum intersection]\label{lem:quorum-intersection}
For any height~\(h\), any two quorums \(Q,Q'\subseteq V_h\) satisfy
\[
  w_h(Q\cap Q')\geq 2q_h-W_h=\iota_h.
\]
Consequently their intersection contains non-Byzantine weight under
Assumption~\ref{ass:height-weights}.
\end{lemma}

\begin{proof}
Since \(w_h(Q\cup Q')\leq W_h\),
\[
 w_h(Q\cap Q')=w_h(Q)+w_h(Q')-w_h(Q\cup Q')
 \geq 2q_h-W_h.
\]
The adversary controls strictly less weight than this lower bound.
\end{proof}

\begin{definition}[Height and current-height target]\label{def:height}
Height is the finality-gadget counter and is separate from slot.  Genesis is treated
as justified and finalized at height \(0\), and the genesis state is at current
height \(1\) with \(T_1=\mathsf{genesis}\).  A justification certificate or timeout certificate
advances height by exactly one.

For \(h>1\), the block whose processing enters height \(h\) is the
\emph{current-height target} \(T_h\), and its slot is the height start slot
\(s_h\).  It is the first block of height~\(h\).  For
height~\(1\), genesis is the target and carries the genesis start slot.  The
target is stored in consensus state until the next height transition.  The
block that becomes \(T_h\) cannot contain a target decision naming itself,
because its root is not yet known when the decision is signed; descendants
can name it.  A checkpoint-sync state carries \(T_h\) even when its trusted
block is in the middle of the height.

The convention \(T_1=\mathsf{genesis}\) applies when the protocol is active
from genesis.  At a fork activation, the first Simplex height starts with an
empty target.  The first eligible post-fork block becomes the target once its
root can be computed; the pre-fork genesis or finalized root is not used.
\end{definition}

\begin{definition}[Blocks and chains]\label{def:block-chain}
A block has a parent, a strictly increasing slot, a deterministic finite set of
signed messages, and a post-state.  Its root is known when its post-state is
recorded.  Blocks form a tree rooted at genesis.
Write \(B\preceq C\) when \(B\) is an ancestor of \(C\), and \(B\sim C\) when
they are compatible.
\end{definition}

\begin{definition}[FG message]\label{def:fg-message}
An FG message by validator \(i\) has a round number, a head root, and two
independent fields:
\[
  m=(i,r,X,d,\phi),
\]
where \(X\) is the head root.
The \emph{height decision} \(d\) is one of
\[
  \mathsf{empty},\qquad
  (\mathsf{timeout},h),\qquad
  (\mathsf{target},h,T).
\]
The finality field \(\phi\) is either \(\bot\) or a commitment
\((h_f,T_f)\) to the latest justified pair in the signer's current chain
state.  An empty height decision has no height and sets neither participation
bit.  Even then, \(X\) may support the next round's pointer and \(\phi\) may
support finalization of the latest justification.
\end{definition}

A target decision is valid on a chain at height \(h\) only when its target is
the stored current-height target \(T=T_h\), and
\(T\) is a strict ancestor of the block containing the vote.  A signed vote for a
different target remains usable as E2 evidence but changes neither the target
bit nor the timeout bit on that chain.  At a nonjustifiable height the state
machine does not justify a target, but a repeated vote for the stored target
still sets the timeout bit.  This allows a validator already locked at that
height to contribute to height advancement without violating E1.  A timeout decision is
valid only at the chain's current height.  A finality commitment to
\((h_f,T_f)\) is counted only while that pair is the chain's latest
unfinalized justification, using \(V_{h_f}\) and \(w_{h_f}\).

\begin{definition}[Slashing rules]\label{def:slashing}
Two signed messages by the same validator are slashable when either condition
below holds.  The two conflicting fields may also occur in the same signed
message.
\begin{description}
  \item[E1:] one carries \(\phi=(h,T)\), while the other's
    height decision at \(h\) is a timeout or targets \(T'\neq T\); or
  \item[E2:] they contain target decisions at the same height for distinct
    targets.
\end{description}
Empty decisions do not conflict with a finality commitment because they make
no height decision.
\end{definition}

E2 makes two target votes at one height slashable.  E1 prevents a validator
that voted to finalize a target from helping a conflicting chain advance past
the same height.

\section{Deterministic finality state machine}\label{sec:state-machine}

The post-state of a block is
\[
  \sigma=(L,s,h,s_h,T_h,\nu_h,\tau,\mu,J,h_j,F,h_F,P).
\]
Here \(L\) is the current block, \(s\) its slot, \(h\) the current height,
\((s_h,T_h)\) the current height's start slot and target, and
\(\nu_h\in\{\mathsf{ordinary},\mathsf{nonjustifiable}\}\) its class.  The
class is set when the chain enters~\(h\) and does not change until the next
height transition.  Finality processing does not change it.  The pair
\((h_F,F)\) is the greatest finalized height and root.  Initially \(h=1\),
\((h_j,J)=(h_F,F)=(0,\mathsf{genesis})\),
\((s_1,T_1)=(0,\mathsf{genesis})\), and \(\nu_1=\mathsf{ordinary}\).  Initially every \(\tau[i]\) and \(\mu[i]\) is
zero and \(P=\varnothing\).  Whenever a block
\(B\) advances height, its post-state sets
\(s_{h+1}=B.\mathsf{slot}\) and \(T_{h+1}=B\).

The target-participation bit \(\tau[i]\) records whether validator~\(i\) has
made a valid decision for the stored current-height target~\(T_h\).  The
timeout bit \(\mu[i]\) records whether it has made either such a target
decision or a valid current-height timeout.  Thus
\[
  Q_{\rm target}(\sigma)=\{i:\tau[i]=1\},
  \qquad
  Q_{\rm timeout}(\sigma)=\{i:\mu[i]=1\}.
\]
The state does not store a target value per validator: the only value that can
set \(\tau[i]\) is the stored target~\(T_h\).  The signed messages are kept so
that E1 and E2 violations can still be proved; an invalid target changes
neither bit.  Different branches may have different current-height targets,
so their participation bits cannot be combined directly.  Each block state is
derived only from its parent and the messages included in that block, as
specified in Section~\ref{sec:store}.

The pair \((h_j,J)\) is the chain's latest justified height and block.  The
set~\(P\) contains validators whose finality fields commit to this pair while
it is the latest justification and has not been finalized.  A newer
justification replaces \((h_j,J)\) and resets~\(P\).  The state does not
keep finality tallies for older justifications.

The signed messages in a block \(B\) are treated as a set.  A valid vote for
the current-height target sets both
\(\tau[i]\) and \(\mu[i]\); a valid timeout sets
only \(\mu[i]\).  A valid finality commitment to the prestate's current
\((h_j,J)\) adds its signer to~\(P\).  Setting bits and adding elements to a
set are independent of message order.  A finality commitment to a
justification created by the same block is not counted.  It must appear in a
descendant whose prestate already contains that justification.  A certificate
received from another branch also does not replace the justified block in the
current branch's state.  Only the justification rule changes \((h_j,J)\)
below.

Whenever the height rules are evaluated, they are applied once in this order.
They may be applied after each block for the safety proof.  In the
round protocol they are applied at the end of the round, after all timely
votes have been processed:
\begin{enumerate}
  \item if \(h_j>h_F\), \(F\preceq J\), and
    \(w_{h_j}(P)\geq q_{h_j}\), set \((h_F,F)\gets(h_j,J)\);
  \item if \(\nu_h=\mathsf{ordinary}\) and
    \(w_h(Q_{\rm target})\geq q_h\), set
    \((h_j,J)\gets(h,T_h)\), reset \(P\), and then advance \(h\), set
    \((s_h,T_h)\) to the current block's slot and root, set the new height's
    class by Definition~\ref{def:nonjustifiable}, and reset \(\tau\) and
    \(\mu\);
  \item otherwise, if \(w_h(Q_{\rm timeout})\geq q_h\), advance \(h\) and
    reset \(\tau\) and \(\mu\), again setting \((s_h,T_h)\) to the current
    block's slot and root and setting the new height's class.  This is the
    timeout branch even when some members of the timeout quorum cast
    current-height target votes.
\end{enumerate}
Finality is processed before a justification or timeout enabled by the same
block.  Since a chain stores one current-height target, all valid target votes
increment the same bitlist.

A target decision therefore has two independent effects: it contributes to
the current-height target tally and sets that validator's timeout bit.  A
validator is counted once in \(Q_{\rm timeout}\), even if the stored messages
later reveal an equivocation.  Justification has priority when the target and
timeout quorums are simultaneously present; otherwise the timeout quorum
advances by timeout.  A validator already locked on the unique target at~\(h\)
need not vote timeout: its existing or repeated target vote already counts
toward the timeout branch.

At the end of a round, finality is checked before justification and timeout.
Thus a full quorum of finality votes for the latest justification is processed
before a new justification can replace it and reset~\(P\).  A partial tally is
discarded when a newer justification replaces it; the protocol finalizes
only the latest justified block.  It needs neither a fork-choice rule for
pending finality nor a per-height finality map.  An empty decision sets neither
participation bit.  Otherwise every validator could submit an empty decision
and advance the height without voting for a target or timeout.

\begin{definition}[Certificates]\label{def:certificates}
A \emph{justification certificate} \(\mathsf{JC}(h,T)\) consists of signed
target votes of weight at least \(q_h\), together with state data showing that
the votes were counted while \(\nu_h=\mathsf{ordinary}\) and \(T_h=T\), and
that the justification rule set \((h_j,J)=(h,T)\).  A target quorum observed
at a nonjustifiable height is not a JC\@.
A \emph{timeout certificate} \(\mathsf{TC}(h)\) is a set of valid
current-height target decisions, timeout decisions, or both, from distinct
validators with weight at least \(q_h\).  When no target quorum is selected first, a
\(\mathsf{TC}(h)\) advances height by timeout.
A \emph{finality certificate} \(\mathsf{FC}(h,T)\) is a
\(\mathsf{JC}(h,T)\) together with finality commitments to \((h,T)\) of weight
at least \(q_h\), collected while \((h,T)\) was the chain state's latest
justification.
Certificates include the ancestry required to validate every nonempty target.
Other validators can verify them, but
receiving a certificate does not change a branch's latest justified block or
authorize a vote on a conflicting chain.
\end{definition}

\begin{lemma}[Target uniqueness]\label{lem:target-uniqueness}
Two conflicting targets cannot both have justification certificates at one
height unless validators of total weight at least~\(\iota_h\) violate E2.
\end{lemma}

\begin{proof}
The signer sets of the two certificates intersect in weight at least
\(\iota_h\) by Lemma~\ref{lem:quorum-intersection}.  Every signer in
the intersection signed both distinct targets at height \(h\).
\end{proof}

\section{Accountable safety}\label{sec:safety}

\begin{lemma}[Height progression]\label{lem:height-progression}
Every height transition increments the height by exactly one and requires a
justification certificate or timeout certificate counted using the validator
weights for that height.
\end{lemma}

\begin{proof}
The justification and timeout rules are the only rules that advance height.  Both
increment it by one, and both use the weights required by
Assumption~\ref{ass:height-weights}.
\end{proof}

\begin{lemma}[Current-height target validity]\label{lem:height-target-freshness}
If a target \(T\) is counted at height \(h\) on a chain ending at \(B\), then
\(T=T_h\preceq B\), and the chain's post-state at \(T\) has current height
\(h\).
\end{lemma}

\begin{proof}
Target validity requires \(T\) to be the stored current-height target and a
strict ancestor of the block containing the vote.  The target is initialized to genesis
for \(h=1\), is set by the block whose transition enters each \(h>1\), and is
replaced by the next height transition, so no height transition occurs between
\(T\) and its use as \(T_h\).
\end{proof}

\begin{lemma}[Target uniqueness on a chain]\label{lem:chain-target-uniqueness}
Each reached height has one current-height target on a given chain.  Suppose a
vote for target~\(T\) is valid on one branch at height~\(h\).  If another
branch at height~\(h\) contains~\(T\), then its current-height target is also
\(T\), and it accepts the vote.  Distinct targets can be accepted only on
conflicting branches.
\end{lemma}

\begin{proof}
Height never decreases and each transition increments it by one.  A chain
therefore has one first block at height~\(h\), which
Definition~\ref{def:height} stores until the next transition.  If another
branch contains that block~\(T\), both branches have the same state at~\(T\)
and entered height~\(h\) there.  Its stored target is consequently also~\(T\).
The target-validity rule then accepts the vote.
\end{proof}

\begin{lemma}[Target participation bits]\label{lem:target-bit-compression}
For a given chain state, the bits \((\tau[i],\mu[i])\) contain all the
information needed for the justification and timeout rules.  Retaining the
signed messages is sufficient to prove E1 and E2 violations.
\end{lemma}

\begin{proof}
By the target-validity rule and Lemma~\ref{lem:chain-target-uniqueness}, every
accepted target vote names~\(T_h\).  Its signer set is therefore
\(\{i:\tau[i]=1\}\).  The set of validators that signed either a valid target
or a valid timeout is \(\{i:\mu[i]=1\}\), with each validator counted once.
These are the only two signer sets used by the height transition.  E1 and E2
are checked from pairs of signed messages, so they do not require target roots
in the participation state.
\end{proof}

\begin{lemma}[Past a finalized height]\label{lem:past-finalized}
Let \(C\) be finalized at height \(h\).  Any block \(B\) whose chain has
advanced beyond \(h\) contains \(C\).  For \(h\geq1\), the only exception
requires E1-slashable validators of total weight at least~\(\iota_h\).
\end{lemma}

\begin{proof}
If \(h=0\), then \(C\) is genesis and the result is immediate.  Assume
\(h\geq1\).
A chain that advances beyond \(h\) contains either a justification certificate
with signer set \(Q_J\), or a timeout certificate with signer set
\(Q_{\rm timeout}\).  The finality certificate for \(C\) has signer set~\(Q_F\).
Either intersection has weight at least \(\iota_h\).  Consider any
non-slashable signer in the relevant intersection.  Its height decision cannot
be a timeout, nor a target different from~\(C\), because either would conflict
by E1 with its finality commitment.  It therefore targeted~\(C\).  Target
validity says that a chain which counted that target vote contains~\(C\).  Thus, if
\(B\) did not contain~\(C\), every member of the quorum intersection would be
E1-slashable, contradicting the bound.
\end{proof}

\begin{lemma}[Finalized blocks form a chain]\label{lem:finalized-chain}
For any finalized pairs \((C,h)\) and \((C',h')\), if \(h\leq h'\) then
\(C\preceq C'\), unless the fault bound in
Assumption~\ref{ass:height-weights} is violated.
\end{lemma}

\begin{proof}
If \(h=0\), then \(C\) is genesis and the result is immediate.  Assume
\(h\geq1\).
If \(h=h'\), intersect a finality certificate for \(C\) with the justification
certificate for \(C'\).  E1 forces every non-slashable intersection signer to
have \(C=C'\).

If \(h<h'\), the chain on which \(C'\) was finalized advanced beyond \(h\), so
Lemma~\ref{lem:past-finalized} implies that this chain contains \(C\).  Both
\(C\) and \(C'\) are therefore on it.  Current-height target validity gives state-height \(h\) at \(C\)
and state-height \(h'\) at \(C'\); because state-height is nondecreasing along a
chain and \(h<h'\), we have \(C\prec C'\).
\end{proof}

\begin{theorem}[Accountable safety]\label{thm:accountable-safety}
Conflicting blocks cannot both finalize unless, for some height \(h\), at least
\(\iota_h\) weight is provably E1- or E2-slashable.
\end{theorem}

\begin{proof}
This follows from Lemmas~\ref{lem:target-uniqueness}
and~\ref{lem:finalized-chain}.
\end{proof}

The proof does not require synchrony.  It is independent of proposer pointers,
the available chain, graded agreement, confirmation, and nonjustifiable
heights.

\section{Store merging}\label{sec:store}

\begin{definition}[Stored data]\label{def:stored-data}
A node stores:
\begin{itemize}
  \item blocks and their post-states;
  \item signed FG and head messages;
  \item pointer proposals, each containing a round and a root;
  \item facts \(\mathsf{timely}(x,I)\), recording that object \(x\) was
    received before the deadline of scheduled step~\(I\);
  \item the result of each scheduled step that has already been completed; and
  \item the local round clock.
\end{itemize}
The first five components are sets.  A timely fact is local metadata, not a
signed network message.  Step~\(I\) uses the objects marked timely for~\(I\)
when its deadline is reached.  Its result is then fixed.  A later message may
be evidence or input to a later step, but it does not change a completed step.
\end{definition}

A block's post-state is obtained from its parent's post-state by processing
only the messages included in that block.  When two stores are merged, each
block is replayed with its own included message batch.  In particular, target,
timeout, and finality bitlists from sibling blocks are never combined.  A
chain accepts only its stored current-height target.  Votes for other targets
can provide slashing evidence, but do not change that chain's state.  Blocks
that conflict with the latest finalized block do not affect the protocol
state.

\begin{definition}[Store merge]\label{def:store-join}
For stores \(S_1,S_2\), their merge \(S_1\sqcup S_2\) contains the union of
their blocks, messages, proposals, timely facts, and completed-step records.
Its clock is the greater of the two clocks.  A scheduled step that neither
copy has completed uses every message marked timely for that step in either
copy.  If these messages contain two distinct pointers signed by the round's
proposer, the step uses neither.  If they contain only one pointer, that
pointer remains available.

A merge does not recompute a completed step.  A second proposal learned after
the pointer deadline is retained as equivocation evidence, but does not
invalidate the pointer already chosen.  If the two copies completed the same
step with different local inputs, the merge retains both histories; it does
not replace them with a single new version of that past step.  A node can use
the merged store for a common continuation after all later steps that depend
on either result have completed.  Conflicting head votes remain stored because
Definition~\ref{def:fresh-root} counts their signer once.  A valid message that
refers to an unknown block remains stored until that block and its ancestors
arrive.
\end{definition}

\begin{theorem}[Order independence of evidence merge]\label{thm:store-merge}
Under the fault bound in Assumption~\ref{ass:height-weights}, and after
ignoring blocks that conflict with the merged finalized root, the union of the stored evidence and the inputs to
future steps is associative, commutative, and idempotent.  Merging \(S_1\) and
\(S_2\) does not require the store from which they were copied.  This statement
does not identify two different completed executions of the same step.
\end{theorem}

\begin{proof}
Set union is associative, commutative, and idempotent; maximum has the same
three properties; and future steps are deterministic functions of their
complete stored input: state, retained results of completed steps, and messages
received before their deadlines.  Completed results are retained as history
rather than recomputed.  By accountable safety, all finalized roots are
compatible under the same bound and therefore have a unique greatest element.
\end{proof}

Head messages are live only for their specified round window.  After merging,
the clock is
\(t=\max(t_1,t_2)\).  An event with expiry \(d\) is live exactly when
\(d>t_1\) and \(d>t_2\).  This remains true if the event is present in only one
input.  An event that has expired in either store is therefore expired in the
merged store, so deleted expired data need not be restored.  A timely fact is
different: \(\mathsf{timely}(x,I)\) remains in the union after the deadline has
passed.

For an available-layer committee selected for a given available vote, the store must retain
the walk head and state from which the committee was chosen, together with the
ordered weighted validator list.  If two stores contain different inputs, the
implementation must either show that they produce the same committee or apply
a deterministic conflict rule.  A cached committee list does not record how
the committee was selected.

Pruning data below a finalized root must retain the finalized state and the
validator weights for relevant heights, the current-height target, the current
height's class and the data used
to determine it, the latest justified pair and evidence that it was justified
at an ordinary height, and any unresolved timely facts.  Validators' signing
histories are local data rather than part of the network store, and are not
pruned by this rule.

\section{Fork-choice state and nonjustifiable heights}\label{sec:fork-choice-state}

Fork choice is computed from the store of Section~\ref{sec:store} as follows.
Let \(\Sigma.\mathcal T\) be the accepted block tree restricted to descendants
of the greatest finalized block~\(\Sigma.F\).  Accepted blocks that conflict
with~\(\Sigma.F\) may remain in the store as evidence, but are not counted by
fork choice.  Among latest justifications in accepted block states whose
blocks descend from~\(\Sigma.F\), let \((\Sigma.h_j,\Sigma.J)\) have the greatest
height, breaking ties by target hash.  Let
\(\Sigma.h_{\max}\) be the greatest state-height of any block in this restricted
tree.  Blocks outside the subtree of \(\Sigma.F\) do not affect the protocol.

\begin{definition}[Viable subtree and cascade root]\label{def:cascade-root}
A block is \emph{viable} in~\(\Sigma\) when it has an accepted leaf descendant
whose state-height is at least \(\Sigma.h_{\max}-1\).  Write
\(\mathcal T'(\Sigma)\) for the viable subtree.  The cascade root is
\[
  R(\Sigma)=
  \begin{cases}
    \Sigma.J,&\Sigma.h_{\max}=\Sigma.h_j+1,\\
    \Sigma.F,&\text{otherwise}.
  \end{cases}
\]
Thus the latest justification roots the walk only while it is exactly one
height behind \(\Sigma.h_{\max}\).  Once a timeout moves a chain one height
farther, the walk starts from the finalized root instead.
Every accepted pointed or fallback anchor below is required to lie in
\(\mathcal T'(\Sigma)\) and to descend from~\(R(\Sigma)\).
\end{definition}

The finalized root is viable because a block that raises \(h_{\max}\) is
accepted only if it descends from the local finalized block.
When the cascade selects~\(J\), the justified target itself is at
state-height~\(h_{\max}-1\), so it is viable as well.  Every descent used below
therefore starts from a viable root.

\begin{lemma}[Finalized lock-in]\label{lem:finalized-lockin}
Under Assumption~\ref{ass:height-weights}, once an honest store has processed the
justification of a finalized block~\(F\), every later accepted block whose
state-height is strictly above~\(F\)'s height descends from~\(F\), and every
later walk output descends from~\(F\).
\end{lemma}

\begin{proof}
Lemma~\ref{lem:past-finalized} gives the first claim.  The store filters later
justifications against its finalized root.  If the cascade root is~\(J\), it
therefore descends~\(F\); if the cascade root is~\(F\), the claim is immediate.
Every anchor and every later step of the walk is constrained to descendants of the
cascade root.
\end{proof}

Fix protocol constants \(K\geq2\) and \(D_{\rm debt}\geq1\).

\begin{definition}[Nonjustifiable height]\label{def:nonjustifiable}
When a chain enters height~\(h\), it first processes finality.  It is in
\emph{finality debt} when the resulting state has
\(h-h_F>D_{\rm debt}\).  The chain sets
\(\nu_h=\mathsf{nonjustifiable}\) exactly when it is in finality debt
and \(K\mid h\); otherwise it sets
\(\nu_h=\mathsf{ordinary}\).  This value is part of the authenticated
consensus state and remains fixed while the chain is at height~\(h\).  At a
nonjustifiable height, the justification branch of the state transition is
disabled.  Once it has confirmed the height, an honest validator with no
earlier vote at that height votes for timeout.  A validator locked on the
chain's current target may repeat that target; the vote sets the timeout bit
but cannot justify.
Every other height is
\emph{ordinary}.

The recovery proof requires a height~\(H\) for which every valid chain reaching
\(H\) that an honest validator can process now or import later has
\(\nu_H=\mathsf{nonjustifiable}\).
This includes hidden historical branches and their authenticated states and
transitions, not only the chains visible during recovery.  Consequently no
valid pre-existing or subsequently assembled \(J(H)\) exists.  While finality
debt continues, one height divisible by~\(K\) occurs every~\(K\) heights, but
seeing one branch enter a nonjustifiable height is not enough.
A concrete protocol must show from authenticated state or transition evidence
that all valid chains entering~\(H\) assign it the same value.
\end{definition}

\begin{lemma}[A nonjustifiable height prevents FG interference]\label{lem:nonjustifiable-fg}
If every valid chain reaching~\(H\) has
\(\nu_H=\mathsf{nonjustifiable}\), no \(J(H)\) can form.  Throughout the
period in which \(\Sigma.h_{\max}=H+1\), every honest
cascade root is its own \(\Sigma.F\); in particular, FG cannot reintroduce a
height-\(H\) justification as the walk root.  If the store already contains,
or later imports, a branch at height~\(H+2\), the only newly eligible
justification root is \(J(H+1)\), not \(J(H)\).
\end{lemma}

\begin{proof}
The state transition never selects the justification branch at~\(H\), even if
compatible locked target repeats reach quorum, so no valid \(J(H)\) exists.
Every chain with \(\Sigma.h_{\max}=H+1\) therefore has greatest justified height
at most \(H-1\).  While \(\Sigma.h_{\max}=H+1\), we have
\(\Sigma.h_{\max}\geq\Sigma.h_j+2\), so the condition for selecting~\(J\) is
not met and Definition~\ref{def:cascade-root} selects~\(\Sigma.F\).  The same
reasoning applies after importing any branch that leaves
\(\Sigma.h_{\max}=H+1\), because a
\(J(H)\) remains impossible.  When \(\Sigma.h_{\max}=H+2\), the cascade can instead select
only \(J(H+1)\).  Lemma~\ref{lem:cascade-compatible} shows that any such root
used during recovery is compatible with the common safe-confirmed chain once
it has been established.
\end{proof}

While finality debt continues, one height divisible by~\(K\) occurs every
\(K\) heights.  The bound \(K\geq2\) ensures that its successor is ordinary;
otherwise every new height could be nonjustifiable and no new target could be
justified.

\section{SG latest messages and fork choice}\label{sec:forkchoice}

The liveness proof considers a period without validator-set changes.  The
following assumption states the required participation and weight bounds.

\begin{assumption}[Validator weights during recovery]\label{ass:recovery-weights}
Throughout recovery, the SG rounds and the FG heights on which they can vote
use one fixed validator set \(V^\star\) with positive integer weights \(w(i)\).
Write
\[
  W=\sum_{i\in V^\star}w(i),
  \qquad
  q=\left\lceil\frac{2W}{3}\right\rceil,
  \qquad
  \iota=2q-W.
\]
For \(S\subseteq V^\star\), write \(w(S)=\sum_{i\in S}w(i)\).
Adversarial weight is \(b<\iota\).  Let~\(V_{\rm resp}\) be a responsive honest
set and write \(W_{\rm resp}=w(V_{\rm resp})\), with \(W_{\rm resp}\geq q\).
In every honest SG view during recovery, let~\(a\) be the total live
latest-message weight outside~\(V_{\rm resp}\), including
both adversarial validators and any temporarily nonresponsive honest validator
whose stale message has not expired.  The exact integer inequality
\(W_{\rm resp}>2a\)
holds in every such view.
Every responsive honest validator remains online through the available vote,
SG pre-vote, FG decision, and message relay in every round used by the recovery
proof.  A reconfiguring
implementation may replace this assumption only with a proof that
the relevant old and new duties preserve the same quorum and intersection
claims.
Each responsive honest validator publishes exactly one SG pre-vote in every
round during recovery, including a round whose FG decision is empty.  Each such
pre-vote remains live through the immediately following round's pointer and
fallback checks and is superseded by that validator's next honest
pre-vote before expiry.  The proof never relies on an expired vote.
\end{assumption}

For the familiar exact \(2/3\) case, \(q=\lceil2W/3\rceil\) and
\(b<W/3\).  All statements below use integer comparisons against~\(q\) or
cross-multiplication; informal fractions are explanatory only.

\begin{definition}[Live latest head vote]\label{def:head-vote}
After the round's Goldfish confirmation step, validator~\(i\) signs one SG
head pre-vote \((i,r,B_i)\).  Its later FG message repeats the same head and
round.  These are two copies of the same logical head vote, not two SG votes.
For equivocation, however, the protocol compares complete signed messages.  A
validator equivocates by signing two distinct SG pre-votes in one round, or two
distinct FG messages in one round, even if the messages name the same head.  It
also equivocates if its SG pre-vote and FG message name different heads.

Before SG scoring, discard every head vote whose head does not descend from
the store's finalized root.  Retain the signed message as evidence, but do not
use it to determine equivocation, the latest message, or the SG denominator.
Among the remaining messages, for expiry window \(W_{\rm head}>0\), the live
latest message of a non-equivocating validator is its unexpired head vote from
the highest round.  A validator that equivocates among these messages
contributes no latest vote.  In an honest view~\(u\), let
\[
 D_u=\sum_{i\in L_u}w(i),
 \qquad
 \operatorname{supp}_u(X)=
 \sum_{\{i\in L_u:X\preceq B_i\}}w(i),
\]
where \(L_u\) is the locally live, unslashed, non-equivocating set.  Support is
subtree-counted, and \(B_i\) denotes validator~\(i\)'s live latest head.
Latest messages are processed when received and need not be included in a
block.
\end{definition}

\begin{definition}[The two SG grades]\label{def:grades}
The \emph{grade-1 fallback anchor} \(G_1(\Sigma,u)\) starts at the cascade
root~\(R(\Sigma)\) and repeatedly enters the unique viable child~\(c\) with
\(3\operatorname{supp}_u(c)\geq2D_u\), stopping when no such child exists.  If
\(D_u=0\), it returns the cascade root.

A block~\(C\) is \emph{grade-0 clear} in~\(u\)'s view when every block
conflicting with~\(C\) has
\(3\operatorname{supp}_u(X)<D_u\).  Equivalently, no conflicting block has
support at least one third of the local live denominator.  When \(D_u=0\),
grade-0 clearance holds by convention.  The \(2/3\) anchor and \(1/3\)
conflict veto are used at every height.
\end{definition}

At most one child of a block can meet the grade-1 threshold because the child
subtree supports sum to at most~\(D_u\).  Grade-0 clearance is a predicate,
not a selected block.

\begin{assumption}[Goldfish properties]\label{ass:goldfish}
Write~\(\Omega\) for the available-layer blocks and votes used for one round's
walk and confirmation.  The function
\(\mathsf{goldfishHead}(A,\Omega)\) deterministically chooses a viable
descendant of anchor~\(A\).  Goldfish confirmation is separate from SG and
safe confirmation.  It requires a strict majority of the participating
committee.  A vote for a descendant supports every ancestor back to~\(A\).
The participating set is fixed after the available vote, and confirmation is
evaluated one slot later.  During recovery, Goldfish satisfies:
\begin{enumerate}
  \item the available-confirmed heads used when honest validators make the SG
    and FG votes considered here are pairwise compatible, including across
    rounds.  The same holds for the confirmations underlying a saved target,
    an active lock, or a justified target used as a cascade root;
  \item if an honest view confirms~\(C\), and at one post-GST available-vote
    step every responsive honest voter starts from an anchor compatible with
    \(C\) and votes for a descendant of~\(C\), then the delayed evaluation
    makes~\(C\) common to all honest views.  Thereafter, Goldfish from any
    compatible honest anchor returns a descendant of~\(C\), for as long as
    the participation and compatible-anchor assumptions continue to hold;
    and
  \item in a slot with an honest proposer, let \(c_{\rm honest}\) be the common
    online honest committee weight and \(c_{\rm adv}\) the committee's total
    adversarial weight.  Then \(c_{\rm honest}>c_{\rm adv}\).
\end{enumerate}
Any block named by a vote or pointer, together with the data needed to validate
it, is available before the vote or pointer is processed.  The compatibility
assumption applies only to the heads listed in the first condition.  If the
recovery argument is extended to later rounds, the same conditions must hold
in those rounds as well.
\end{assumption}

The third condition holds separately in each view.  Every honest participating
set contains all common honest committee weight~\(c_{\rm honest}\) and may
contain an arbitrary adversarial subset of weight~\(a\leq c_{\rm adv}\).  Hence
\[
  c_{\rm honest}>c_{\rm adv}\geq a
  \quad\Longrightarrow\quad
  c_{\rm honest}>\frac{c_{\rm honest}+a}{2}.
\]
The honest proposal therefore has a strict local majority in every honest
view.  The participating sets and their denominators need not be identical.
The first condition above provides the only compatibility assumption between
different views.

\begin{definition}[Anchored walk]\label{def:walk}
At a round's available-vote action, validator~\(u\) first evaluates the
round-start proposal's pointed root as in
Definition~\ref{def:fresh-root}.  If the root~\(P\) is accepted,
\(R(\Sigma)\preceq P\), and \(P\) is viable, set the anchor~\(A=P\).
Otherwise set \(A=G_1(\Sigma,u)\).  From~\(A\), follow
\(\mathsf{goldfishHead}(A,\Omega)\) inside the viable subtree.  If the result
has state-height below \(\Sigma.h_{\max}-1\), continue through viable children
using the same previous-slot participating-vote plurality, with deterministic
ties, until that height bound is met.  The returned block is
\(\mathsf{walk}(\Sigma,\Omega,u)\).
\end{definition}

The last step affects only the height of the returned block, not its
confirmation status.  By definition, a viable child exists until the height
bound is reached.

\begin{definition}[Safe confirmation]\label{def:safe-confirm}
Fix an honest view~\(u\), its derived store~\(\Sigma\), cascade root
\(R=R(\Sigma)\), and Goldfish available-confirmed head~\(A_u\).  The cascade
root~\(R\) is always safe-confirmed.  If \(A_u\) does not descend from~\(R\),
the safe-confirmed head is~\(R\).  Otherwise it is the deepest grade-0-clear
block between~\(R\) and~\(A_u\), or~\(R\) if no strict descendant qualifies.
Every block between~\(R\) and the returned block is also safe-confirmed.  This
rule is used only by FG; it does not change the Goldfish available-confirmed
head.
\end{definition}

The prefix property follows because every branch conflicting with an ancestor
of a grade-0-clear block also conflicts with that block.

\begin{assumption}[Post-GST delivery and an honest proposer]\label{ass:honest-proposer}
After GST, honest blocks, available votes, SG pre-votes, FG messages, and the
blocks needed to validate them propagate within~\(\Delta\).  If the first
proposer of a round is honest, then:
\begin{enumerate}
  \item the preceding round's responsive honest SG pre-votes reach the proposer
    before it chooses its root.  The proposal, the block and vote state fixed
    by that proposal, and the pointed root reach every honest validator
    at least one delivery bound before the available-vote action;
  \item the proposer re-sends its supporting votes, so its pointed root has
    credited support at least~\(q\) in every honest view by
    Lemma~\ref{lem:fresh-root-sync};
  \item every honest validator evaluates pointer ancestry and viability using
    the same state.  If the pointed root passes, all use it; if it fails,
    all use the same grade-1 fallback.  Thus every honest available voter uses
    one common viable anchor and votes for the proposal block or a descendant
    on its walk chain; and
  \item after delayed confirmation, every honest validator chooses an SG head
    on a descendant of the same block and obtains the same height and
    current-height target.  No validator switches to unrelated local state
    between the available vote and the SG vote.
\end{enumerate}
An implementation that spreads these duties across slots must ensure the same
conditions.  They do not assume a height advance, justification, or finality.
\end{assumption}

\section{Fresh-quorum pointers}\label{sec:pointer}

A pointer is a synchronization hint, not a certificate carried by the block.
Its support is evaluated locally from the immediately preceding round's signed
head messages.

\begin{definition}[Credited support and fresh root]\label{def:fresh-root}
Fix honest view~\(u\), support round~\(r\), and known block~\(P\).  Let
\(V_u^r(P)\) be the validators from which~\(u\) has a valid round-\(r\) head
vote descending from~\(P\), and let \(E_u^r\) be the validators for which~\(u\)
has observed an equivocation under Definition~\ref{def:head-vote}.  Two
distinct same-round FG messages count even if their head roots agree.  Both
sets use only messages received by the
deadline for this round's available vote, as recorded by the store's
\(\mathsf{timely}\) facts.  A later message may support a later round but does
not change this decision.  Define
\[
  \operatorname{credit}_u^r(P)
   =w\bigl(V_u^r(P)\cup E_u^r\bigr).
\]
Each validator is counted once.  An equivocator is credited to~\(P\) even when
none of the two copies in~\(u\)'s view descends from~\(P\).

The first proposer of round~\(r+1\) may point to one block root~\(P\).  View~\(u\)
accepts it as a \emph{fresh root} at its available-vote step exactly when
\(P\) and its ancestry are available,
\(\operatorname{credit}_u^r(P)\geq q\), the proposer has not signed a distinct
pointer for the same round, and the walk checks in
Definition~\ref{def:walk} (cascade ancestry and viability) pass.  The proposal
carries only~\(P\), never the votes or an aggregate.
\end{definition}

Before the deadline, credited support for a fixed~\(P\) can only increase.  A
supporting vote adds its signer; another vote from the same signer either leaves
that support unchanged or proves an equivocation, which is also credited.  At
the deadline, the \(\mathsf{timely}\) records determine which messages are
used.  Since~\(W\) is fixed, the proposal does not need to carry or negotiate a
separate quorum snapshot.

\begin{lemma}[Fresh-support intersection]\label{lem:pointer-intersection}
If two conflicting roots~\(P,P'\) each have credited support at least~\(q\) in
possibly different honest views for the same support round, at least
\(\iota\) weight actually equivocated in that round.  In
particular, two conflicting fresh roots cannot both be accepted under
\(b<\iota\).
\end{lemma}

\begin{proof}
The two credited sets are quorums, so their intersection has weight at least
\(\iota\).  A signer explicitly present in either local
equivocation set has already equivocated.  Otherwise the two views contain one
signed head from that signer descending~\(P\) and one descending~\(P'\).
Because the roots conflict, the heads are distinct.  Thus every intersection
signer actually equivocated, regardless of which copies either verifier saw.
\end{proof}

\begin{lemma}[An honest proposer's fresh root is accepted]\label{lem:fresh-root-sync}
After GST, suppose the first proposer of round~\(r+1\) is honest and points to
a root~\(P\) for which it saw round-\(r\) supporting votes of weight at
least~\(q\), without counting equivocations.  Suppose it re-sends those votes
at least~\(\Delta\) before every honest available-vote action.  By each such
action, \(P\) has credited support at least~\(q\) in that validator's view.
\end{lemma}

\begin{proof}
For each signer counted directly by the proposer, the proposer re-sends the
valid supporting head message immediately, though the proposal itself still
carries only the root.  The post-GST delivery schedule puts that message in
every honest view before the available-vote action.  If a receiver instead has
two different same-round messages from that signer, it counts the signer as an
equivocator.  Thus every signer counted by the proposer is counted either as
support or as an equivocation by every receiver.  Every receiver therefore
credits at least~\(q\) weight to~\(P\).
\end{proof}

This is precisely the benefit-of-the-doubt case: the proposer may have seen
validator~\(i\) vote for~\(P\), while another validator has only two other
equivocating copies.  The latter still credits~\(i\) to~\(P\).  Synchrony or
equivocation is enough; no receiver attempts to reconstruct the proposer's
exact vote set.

\begin{lemma}[A directly supported root exists]\label{lem:fresh-root-exists}
After one network delay for honest SG votes, the highest common ancestor of all
responsive honest round-\(r\) heads has direct support of weight at least~\(q\)
and is known to an honest first proposer of round~\(r+1\).
\end{lemma}

\begin{proof}
Each responsive honest validator emits one round-\(r\) pre-vote and their total
weight is at least~\(q\).  Their highest common ancestor lies on every such
head.  Delivery makes all these messages and the ancestor available to the
honest proposer, which therefore sees direct quorum support.
\end{proof}

An honest proposer chooses the deepest directly quorum-supported root that
descends from the cascade root and passes the viability check in the state used for its proposal,
with a deterministic hash tie-break; if none passes, it omits the pointer and
the fallback is used.  It broadcasts one proposal.  A
Byzantine proposer may equivocate.  If both proposals are in a validator's
view before its pointer deadline, that validator uses neither pointer.  If the second
proposal arrives later, it is kept as equivocation evidence but does not alter
the anchor already chosen for that round.  This can delay liveness but cannot
create two conflicting accepted fresh roots under \(b<\iota\).  There is
a new pointer opportunity every round, even when the height has not advanced.
A pointer applies to one round, not to an entire height.

\begin{lemma}[An honest proposer gives validators the same fork-choice input]\label{lem:pointer-alignment}
Under Assumptions~\ref{ass:recovery-weights}
and~\ref{ass:honest-proposer}, an honest first-slot proposer makes every responsive
honest validator use the same viable anchor at its available-vote action: the
pointed root if its common checks pass, otherwise the common grade-1 fallback.
At the later SG action their walk outputs lie on one chain and have one
state-height~\(k\) and one current-height target~\(T\), even if their pre-round walks
were split between \(h_{\max}-1\) and \(h_{\max}\).
\end{lemma}

\begin{proof}
Lemma~\ref{lem:fresh-root-exists} supplies a directly supported root and
Lemma~\ref{lem:fresh-root-sync} makes its credited support reach~\(q\)
everywhere.  Those two facts alone do not prove cascade-root ancestry and
viability.  Assumption~\ref{ass:honest-proposer} gives them the same state:
the proposed root either passes everywhere or fails everywhere, in which case
the same fallback is used.  It then makes all honest available voters follow
that common anchor's walk chain.  After confirmation, all SG walks remain on
that chain.
The height and current-height target are read from those walk states rather
than from a receiver's pre-round height, so a prior one-height split does not
cause pointer rejection.  The remaining conditions of
Assumption~\ref{ass:honest-proposer} ensure that proposal processing and
Goldfish confirmation give the same height and current-height target.
\end{proof}

\section{Round schedule and state processing}\label{sec:rounds}

FG time is grouped into rounds of 32 consecutive slots.  The paper does not
spread FG duties across those slots; it models one quorum-contributing FG
action per validator after the proposal has received Goldfish available votes
and been confirmed.  The strict order is
\[
\begin{array}{c}
\text{first-slot proposal and pointer selection}
\longrightarrow \text{Goldfish available vote}
\\[-1mm]
\longrightarrow \text{delayed Goldfish confirmation}
\longrightarrow \text{signed SG head pre-vote}
\\[-1mm]
\xrightarrow{\ \Delta\ }
\text{FG checkpoint/finality action}
\\[-1mm]
\longrightarrow \text{process the round's outcome}.
\end{array}
\]
Before a step reads state, it processes every message whose deadline has
passed.
The proposal and available vote occupy the first available-layer opportunity;
the time-shifted confirmation is evaluated one slot later; the SG pre-vote is
then delivered for one network delay~\(\Delta\) before the FG action.  The remaining
slots are available for inclusion and processing.

\begin{lemma}[Delivery before the next step]\label{lem:delivery-before-action}
After GST, every message sent by an honest validator in one step is received
by every honest validator before the next step.  In particular, all honest
available votes arrive before delayed confirmation, all honest SG pre-votes
arrive before the FG action, and all honest FG messages arrive within one
further~\(\Delta\).
\end{lemma}

\begin{proof}
Successive steps are separated by at least the post-GST delay bound~\(\Delta\),
and messages that arrive on time are processed before the next step.  Apply this once to each
arrow in the displayed schedule.  A Byzantine message may arrive before a
deadline in some views and after it in others; the lemma applies only to
honestly sent messages.
\end{proof}

\begin{definition}[Round close]\label{def:round-close}
At the end of a round, the protocol processes all timely votes included by the
specified deadline, then evaluates the state transition once.  It
checks finality first, current-height target justification second, and timeout
third.  It does not evaluate any branch earlier using only some of the round's
timely votes.
\end{definition}

In every round counted below, responsive honest validators of total weight at
least~\(q\) complete the available, SG, and FG steps, and their votes are
included by the deadline for that round close.  If a round advances the
height, before the next SG and FG votes every responsive honest validator uses
the same resulting state and has safe-confirmed a block carrying it.  The
first counted round has an honest first-slot proposer whose
proposal is common before the available-vote action.  Later counted rounds do
not require honest proposers once the new state and its confirmed chain are
common.

This uses the state transition of Section~\ref{sec:state-machine}.  Safety does
not depend on whether it runs immediately or at the end of the round.  The
liveness bounds assume end-of-round processing so that all timely votes are
included before the target branch is checked.

\begin{lemma}[Latest-message concentration]\label{lem:head-concentration}
In the exact \(2/3\), less-than-\(1/3\) recovery regime, suppose all responsive
honest validators publish their next SG pre-vote on descendants of one
chain~\(C\).  One delivery bound later, every honest latest message from the
responsive set lies on~\(C\)'s chain in every honest view.  Every conflicting
branch then has support below one third of the local live denominator, and
remains below it while those honest latest messages stay on the chain.
\end{lemma}

\begin{proof}
Let \(W_{\rm resp}\) be the common responsive honest weight and \(a\) the
weight of live latest messages outside that set in one view.  The recovery assumption
gives \(W_{\rm resp}>2a\).  A conflicting branch has support at most~\(a\), while
\(D=W_{\rm resp}+a\); hence \(3a<W_{\rm resp}+a=D\).  Later honest messages only supersede
same-validator stale messages with another descendant of~\(C\).
\end{proof}

A deployment using a nonstandard rounded threshold must check
\(W_{\rm resp}>2a\) directly
with integer weights; it must not infer the inequality from a rounded adjusted
denominator.

\section{Validator voting rule}\label{sec:signer}

For each height, an honest validator records
\[
  \mathsf{firstTarget}[h],\qquad
  \mathsf{timedOut}[h],\qquad
  \mathsf{lockTarget}[h].
\]
It also records its SG and FG messages for every round.  These records are
written before the validator releases a signature.  After a restart, the
validator does not sign until it has recovered them.  A lock retains the
corresponding justification certificate and ancestry;
\(\mathsf{firstTarget}[h]\) is written at most once, while
\(\mathsf{timedOut}[h]\) can change only from false to true.  Once validators
again use the same chain and have confirmed its current-height target, a
validator repeats its saved target if it is on that chain or, if it is not
locked, votes for timeout.

An honest validator makes a target or timeout vote at height~\(k\) only from
an SG walk
whose stored post-state has reached~\(k\).  In particular, signing at height
\(k\) therefore proves that the signer previously reached that height.  The
record is retained across restarts and checkpoint sync.

\begin{definition}[Current-height target]\label{def:current-height-target}
Let \(W_r\) be the walk output chosen by the round's SG pre-vote and let
\(k=\sigma[W_r].h\).  The round's target~\(T\) is the first block on
\(W_r\)'s chain whose post-state is at height~\(k\), equivalently that chain's
authenticated current-height target~\(T_k\).  The current-height decision is about~\((k,T)\),
not necessarily about the pointed root itself.  Validators first choose their
walk, then derive its height and target.
\end{definition}

\begin{definition}[Finality vote]\label{def:finality-vote}
Let \((h_j,J)\) and \((h_F,F)\) be the latest justified and finalized pairs in
the chain state used by the FG action.  An honest validator sets
\(\phi=(h_j,J)\) exactly when \(h_j>h_F\), \(F\preceq J\), and:
\begin{enumerate}
  \item it previously signed target \(J\) at \(h_j\);
  \item it has never signed timeout at \(h_j\);
  \item it knows the justification certificate for \(J\); and
  \item its lock entry is empty or already equals \(J\).
\end{enumerate}
Before signing, it records \(\mathsf{lockTarget}[h_j]\gets J\) and retains the
certificate.  Otherwise it sets \(\phi=\bot\).  A newer justification becomes the only eligible finality
commitment and resets the previous justification's partial finality tally.
The rule finalizes only the
latest justified block.
\end{definition}

\begin{definition}[FG voting rule]\label{def:fg-rule}
At the FG action, let \((k,T)\) be derived from the SG walk chosen earlier in
the round, and write~\(C_i\) for validator~\(i\)'s safe-confirmed head.  After independently
choosing its finality vote, validator~\(i\) chooses its current-height decision
in this order:
\begin{enumerate}
  \item If \(\mathsf{lockTarget}[k]=T_L\neq\bot\), repeat target~\(T_L\) iff
    \(T_L=T\) and \(C_i\succeq T\); otherwise emit empty.  The vote may be
    included only on a descendant of~\(T\).  At a nonjustifiable height the repeated
    target sets the timeout bit but cannot justify.  A locked
    validator never signs timeout at~\(k\).
  \item If \(\mathsf{timedOut}[k]\) is true, repeat timeout iff
    \(\sigma[C_i].h\geq k\); otherwise emit empty.
  \item If \(\mathsf{firstTarget}[k]=T_0\neq\bot\) without a lock, repeat
    target~\(T_0\) iff \(T_0=T\) and \(C_i\succeq T\).  Such a vote may be
    included only on a descendant of~\(T\).  Otherwise sign timeout iff
    \(\sigma[C_i].h\geq k\); otherwise emit empty.  Thus a validator repeats a
    saved target that is on the current chain, and switches to timeout if that
    target is not on the chain.  At a nonjustifiable height an exact repeat
    sets the timeout bit but cannot justify.  Signing the timeout
    permanently forfeits finality eligibility for the old target.
  \item With no history at a nonjustifiable height, sign timeout iff
    \(\sigma[C_i].h\geq k\); otherwise emit empty.
  \item With no history at an ordinary height, sign target~\(T\) iff
    \(C_i\succeq T\); otherwise sign timeout iff
    \(\sigma[C_i].h\geq k\); otherwise emit empty.
\end{enumerate}
A new target writes \(\mathsf{firstTarget}[k]=T\) before signing; a timeout sets
\(\mathsf{timedOut}[k]\).  Every message repeats the already-signed SG
head~\(W_r\).  An empty decision has no height and contributes to neither a
target nor a timeout quorum, but still carries that head and the independently
computed finality field.
\end{definition}

Timeouts require confirmation of the height, but do not require an additional
grade-0 check.  Once a validator safely confirms the current-height target, it
contributes either a target or a timeout.  If it has not confirmed that height,
it emits an empty decision.  This does not exempt it from a stall penalty.

\begin{lemma}[Honest votes are not slashable]\label{lem:signer-safety}
The validator of Definition~\ref{def:fg-rule} never violates E1 or E2.
\end{lemma}

\begin{proof}
The first-target entry can be set only once, so the validator never signs two
distinct targets.  It may sign timeout after a target that cannot be repeated on the
current confirmed chain only while no lock exists; after doing so it can never
make a finality vote at that height.  Once a finality
commitment creates a lock, the validator either repeats the exact locked target or
emits empty.  It never signs a conflicting target or timeout after the
commitment.
\end{proof}

\begin{definition}[Using earlier votes and certificates]\label{def:earlier-votes}
While a height can affect an honest fork choice, every honest validator
re-sends:
\begin{enumerate}
  \item every saved target or timeout vote that is still valid on a branch an
    honest validator may follow, including one from a past local height that
    is still current on a competing chain;
  \item the blocks, ancestry, and signed evidence needed to validate the
    branch's latest justified/finalized state, plus any justification
    certificate retained by an active lock; and
  \item every live SG head message needed for the latest-message and
    previous-round pointer-support checks.
\end{enumerate}
Votes whose referenced blocks are not yet known are queued and processed once
those blocks arrive.  A signature
does not affect a chain's consensus state until it is included in a block
whose transition accepts it; a certificate by itself does not replace the
branch's latest justified pair.  The validator may discard an entry only once
every honest validator has finalized the same root and no competing branch at
or below that height can affect protocol choices.  Advancing its local
height is not enough.
\end{definition}

\begin{lemma}[Locked validators still contribute to height advance]\label{lem:lock-alignment}
Suppose that, during recovery, all responsive honest validators use one chain
at height~\(k\), derive its current-height target~\(T\), and safely confirm
that target before voting.  If an honest validator is locked at~\(k\), then
either the chain has already advanced beyond~\(k\), or the validator's saved
target is~\(T\) and its earlier or repeated target vote contributes to the
timeout quorum.
\end{lemma}

\begin{proof}
A lock at~\(k\) was created only after the validator voted for~\(T_L\) and
learned \(\mathsf{JC}(k,T_L)\).  If that target quorum has already been
processed on the current chain, the height has advanced.  Otherwise the common
walk is at height~\(k\) and its safely confirmed target is~\(T\).

If \(T\) is the cascade root, it is genesis or a justified or finalized target
at the same height.  Target uniqueness and finalized-chain safety, applied to
the retained \(\mathsf{JC}(k,T_L)\), give \(T_L=T\) under \(b<\iota\).
Otherwise, Goldfish confirmed a descendant of~\(T\).  The retained JC for
\(T_L\) contains an honest signer who also Goldfish-confirmed a descendant of
\(T_L\) before voting.  Assumption~\ref{ass:goldfish} makes these confirmations
compatible.  Thus \(T_L\) and~\(T\) lie on one chain; because both are its
current-height target at height~\(k\), they are equal.  The original
target vote contributes to the timeout quorum if it has already been included;
otherwise the validator repeats~\(T\).  It never needs to vote for timeout.
\end{proof}

\begin{remark}[Leak-fairness consequence]\label{rem:leak-fairness}
The target bit does not assume that competing chains have one shared target.
Each chain accepts only its own current-height target.  Once an honest proposer
makes all responsive honest validators use the same chain, an unlocked
validator repeats a saved target on that chain.  If its saved target is not on
the chain, it may vote for timeout because it has made no finality commitment.
A locked validator cannot vote for timeout, but
Lemma~\ref{lem:lock-alignment} shows that its justified target lies on the
chain now used by honest validators and equals its current-height target.  Its
target therefore sets the timeout bit just as a timeout vote would.  The target
bit therefore preserves every honest validator's contribution to the timeout
quorum.
\end{remark}

\section{Healing and finality liveness}\label{sec:healing}

We first show that two height advances from a nonjustifiable height stabilize
the available chain.  We then give the additional conditions under which a
new justification forms and finality resumes.  Stabilizing the available
chain does not by itself imply finality.

In this section, a \emph{shared store}~\(\Sigma\) means that every responsive
honest validator uses~\(\Sigma\) as input to the next protocol step.  A
validator may retain additional evidence that cannot affect that step.

\begin{assumption}[Message delivery and proposer fairness]\label{ass:recovery-delivery}
After GST, every message sent by an honest validator in one step arrives at
every honest validator before the next step and before the deadline at which
it is used.  Honest current-height and finality votes are included on the
common chain by the applicable round-close deadline.  There is a future honest
first-slot proposer.  After each common height transition, another honest
proposer eventually proposes a descendant of the new state.

The numerical round bounds require this timely processing in every counted
round, with the additional conditions stated after
Definition~\ref{def:round-close}.  Without continuous timely processing, the
same argument gives eventual height advance by waiting for an honest proposal
after each transition.  Eventual finality then requires a later ordinary
height above every height previously reached or signed at by an honest
validator, reached before any honest validator votes at that height.  Its
target votes and the following finality votes must also be included by their
respective deadlines.  Recurring honest proposers alone do not guarantee such
a height, because votes may be cast there before the chain becomes common.
\end{assumption}

\begin{definition}[Common target before FG voting]\label{def:common-target}
For round~\(r\), let \(W_i^r\) be honest validator~\(i\)'s SG walk output,
\(k_i^r=\sigma[W_i^r].h\), and \(T_i^r\) its current-height target.  The round
has a \emph{common target \((k,T)\)} when every responsive honest validator
has \(k_i^r=k\), \(T_i^r=T\), and all \(W_i^r\) lie on one chain.  The target
is \emph{confirmed before FG voting} when every responsive honest validator's
safe-confirmed head lies on that chain at height at least~\(k\).  Every such
head then descends~\(T\).
\end{definition}

\begin{lemma}[An honest proposer gives validators a common confirmed target]\label{lem:honest-proposer-confirmation}
Under Assumptions~\ref{ass:goldfish} and~\ref{ass:honest-proposer}, a post-GST
round with an honest first-slot proposer either observes an already completed
height transition or has a common target that is confirmed before the round's
FG action.
\end{lemma}

\begin{proof}
Lemma~\ref{lem:pointer-alignment} gives one chain, height~\(k\), and
current-height target~\(T\).  Let~\(B_r\) be the honest proposal block.  The
state used for the proposal includes every signed message received by the
applicable deadline and the preceding round's state transition.  Because
the height rules run only at round close, no further height transition occurs
between this state and the round's SG action.  If
this processing already moves the chain beyond~\(k\), the first alternative
holds.  Otherwise \(T\preceq B_r\), and every honest available voter votes for
\(B_r\) or a descendant.  Write~\(a\leq c_{\rm adv}\) for the adversarial
weight in one honest view.  Assumption~\ref{ass:goldfish} gives
\(c_{\rm honest}>c_{\rm adv}\), hence
\(c_{\rm honest}>(c_{\rm honest}+a)/2\).  Delayed evaluation therefore confirms~\(B_r\), and its
ancestor~\(T\), in every honest view without requiring equal denominators.

Every honest validator then publishes its SG pre-vote on that chain.
One delivery bound later Lemma~\ref{lem:head-concentration} makes every
conflicting branch strictly sub-\(1/3\), so~\(T\) is grade-0 clear as well as
available-confirmed.  Thus every safe-confirmed head has reached the height
before the FG action.
\end{proof}

\begin{lemma}[A common confirmed target advances the height]\label{lem:aligned-height-advance}
Consider a round that satisfies Assumption~\ref{ass:goldfish}.  Suppose it has a
common target~\((k,T)\), that target is confirmed before FG voting, and the
responsive honest votes together with all unresolved older votes that they
retransmit are processed before the round closes.  The responsive honest
validators have total weight at least~\(q\) and set their timeout bits at~\(k\).
The round advances the height exactly
once: by justification if one target has weight~\(q\), and otherwise by
timeout.  At an ordinary height, if no honest validator previously timed out
and no honest validator voted for a target other than~\(T\), then~\(T\) is
justified.
\end{lemma}

\begin{proof}
Partition the responsive honest validators by their previous votes at
height~\(k\).  A previously processed target or timeout already set the
validator's timeout bit.
Validators that previously timed out repeat their timeout.  An unlocked old
target that lies on the common chain equals its current-height target~\(T\) by
Lemma~\ref{lem:chain-target-uniqueness} and is repeated; one that does not lie on the
current confirmed chain votes for timeout.  For a lock,
Lemma~\ref{lem:lock-alignment} says either the height has
already advanced or the saved target is~\(T\), and its old or repeated vote
sets the timeout bit.  A validator that has not yet voted at~\(k\) targets~\(T\) at an
ordinary height and times out at a nonjustifiable height.  Confirmation
of~\(T\) makes every validator satisfy the corresponding condition in
Definition~\ref{def:fg-rule}, so none emits an empty vote.

Every responsive honest validator has therefore set its timeout bit, and their
total weight is at least~\(q\).  By assumption, the old votes and the newly
repeated targets or timeouts are processed before the same
round closes.  The state transition first checks for a target quorum and then
for a timeout quorum, and increments the height only once.  Under the final
conditions of the lemma, every previous honest target is~\(T\), every honest
validator that has not yet voted targets~\(T\), and no honest validator has
timed out.  Target support therefore reaches~\(q\).
\end{proof}

\begin{corollary}[One honest-proposer round advances one height]\label{cor:honest-round-advance}
A post-GST round satisfying Assumption~\ref{ass:goldfish}, with an honest
first-slot proposer and the message inclusion of
Assumption~\ref{ass:recovery-delivery}, advances the common height when the round
closes, unless a certificate has already advanced it.  The advance need not
be a justification: if validators previously cast both target and timeout
votes at that height, only the timeout quorum may reach~\(q\).
\end{corollary}

\begin{proof}
Combine Lemmas~\ref{lem:honest-proposer-confirmation}
and~\ref{lem:aligned-height-advance}.
\end{proof}

\begin{corollary}[At most one round is needed to reach \(h_{\max}\)]\label{cor:height-catchup}
Suppose the shared store after processing a proposal has
\(\Sigma.h_{\max}=H\), so the
height-filter postcondition permits an honest walk only at~\(H-1\) or~\(H\).
An honest-proposer round gives all responsive honest validators a common
height and advances it.  If that height is~\(H-1\), one round brings the common
chain to~\(H\).  The next timely round, or the next honest proposal after the
new state is common, acts at~\(H\) unless a certificate has already advanced
it.  Validators therefore need not all be voting at~\(H\) before the first
honest proposal.
\end{corollary}

\begin{proof}
Lemma~\ref{lem:pointer-alignment} derives one height from the common
post-proposal walk.  Corollary~\ref{cor:honest-round-advance} advances that
height exactly once.  The next timely round, or the next honest proposal,
starts from the resulting common state.
\end{proof}

\begin{lemma}[A timeout quorum contains an honest confirmation]\label{lem:timeout-honest-confirmation}
Under Assumption~\ref{ass:height-weights}, whenever a timeout quorum advances
height~\(h\), at least one honest validator in that quorum safe-confirmed a
block at height~\(h\) before voting.  Empty votes do not count.  Since the same
timeout vote may be included on sibling chains at the same height,
the timeout quorum alone does not identify the current-height target of the
chain that includes it.
\end{lemma}

\begin{proof}
The timeout certificate has weight at least~\(q_h\), while adversarial weight is
strictly below the quorum-intersection bound and in particular below~\(q_h\).
It therefore contains an honest decision.  Every honest target decision
requires a safe-confirmed head descending the current-height target; every honest
timeout requires the safe-confirmed head to have reached height~\(h\)
in the signing view.  An empty vote has no height and is excluded from the
timeout quorum.  A timeout therefore proves that an honest signer confirmed some
block at height~\(h\), but not which target belongs to the chain that later
includes the timeout.
\end{proof}

\begin{lemma}[Votes advancing a common height share a confirmed target]\label{lem:common-target-confirmation}
Suppose a round satisfying Assumption~\ref{ass:goldfish} has common
target~\((k,T)\), that~\(T\) is confirmed before FG voting, and responsive
honest votes of weight at least~\(q\) are processed before the round closes.
Each of these votes is a target or timeout vote cast with a safe-confirmed
head descending from~\(T\).  Thus every responsive honest signer has
safe-confirmed~\(T\), and~\(T\) remains a common prefix whether the height
advances by justification or timeout.
\end{lemma}

\begin{proof}
Every responsive honest safe-confirmed head lies on the common chain at height
at least~\(k\), hence above its authenticated current-height target~\(T\).
For a lock, Lemma~\ref{lem:lock-alignment} gives either an already
completed transition or equality with~\(T\); in the nonadvanced case
Definition~\ref{def:fg-rule} makes it repeat that same~\(T\).  A validator that
previously timed out, or whose saved target is incompatible, signs a new
timeout only after the safe confirmation in the premise.  A validator with a
compatible saved target repeats it, and a validator with no previous vote
chooses the common target or timeout according to the height class.  Thus each
current message sets the
timeout bit while its signer safe-confirms~\(T\).  Their recorded weight is at
least~\(q\), so they themselves form a timeout quorum even if older target or
timeout votes are also present.  Whether the target or timeout branch advances the
height does not change the common confirmed prefix established by these votes.
\end{proof}

\begin{lemma}[Later cascade roots are compatible with a safe-confirmed block]\label{lem:cascade-compatible}
Let~\(C\) be safe-confirmed during an interval in which
Assumption~\ref{ass:goldfish} holds.  Under \(b<\iota\), every cascade root
selected at that or a later action in the same interval is compatible
with~\(C\).
\end{lemma}

\begin{proof}
Every non-genesis justified block used as a cascade root has a JC containing an
honest target vote.  Before that vote, the validator Goldfish-confirmed a
descendant of the target.  Every non-genesis finalized block was previously
justified and has the same property.

If~\(C\) lies strictly above its action's cascade floor, the current
available-confirmed head confirms~\(C\).  If~\(C\) is the floor, then it is
genesis or a justified or finalized root with such a prior confirmation.
Clause~1 of Assumption~\ref{ass:goldfish} makes these confirmations
pairwise compatible, so their ancestors are compatible with~\(C\).  This
applies to every justification selected during the interval.  Since a cascade
root is either~\(F\) or~\(J\), later updates to either value preserve the
result.
\end{proof}

\begin{lemma}[Delivered SG votes exclude later conflicting anchors]\label{lem:no-conflicting-anchor}
Suppose~\(C\) is on the common safe-confirmed chain, all responsive honest
available voters perform one post-GST available-vote action on descendants
of~\(C\) from compatible anchors, and all responsive honest validators then
publish that round's SG pre-votes for descendants of~\(C\).  After their
delivery, no anchor adopted by an honest validator in that or a later round
in which Assumption~\ref{ass:goldfish} holds conflicts with~\(C\).
\end{lemma}

\begin{proof}
The walks descend both their fixed anchors and~\(C\).  Their anchors are
therefore compatible.  Lemma~\ref{lem:cascade-compatible} makes every later
cascade root compatible with~\(C\).  After the SG votes are delivered,
Lemma~\ref{lem:head-concentration} leaves every conflicting branch below the
grade-0 threshold.  A grade-1 fallback descent from a compatible root cannot
enter such a branch because every step needs at least two-thirds support.

For a pointed anchor, let~\(P_C\) be the highest common ancestor of the
preceding round's responsive honest pre-vote heads.  It descends from~\(C\) and
has direct support at least~\(q\).  If an honest view credited a conflicting
pointed root~\(P\) with weight~\(q\), intersect its credited signer set with
the direct signer set of~\(P_C\).  The intersection argument of
Lemma~\ref{lem:pointer-intersection} forces at
least \(\iota\) actual equivocation; \(P_C\) need not itself have appeared in a
proposal.  Thus no conflicting pointer is accepted under \(b<\iota\).
The persistence condition in Assumption~\ref{ass:goldfish} keeps the resulting
walks on~\(C\)'s chain, so later honest pre-votes do the same.  The result
follows by induction over later rounds in which the assumption holds.
\end{proof}

\begin{theorem}[Two advances from a nonjustifiable height stabilize the available chain]\label{thm:available-healing}
Let every valid branch reaching~\(H\), including one imported after recovery
begins, have \(\nu_H=\mathsf{nonjustifiable}\), and suppose
either:
\begin{enumerate}[label=(\alph*)]
  \item a round has a common target at~\(H\), that target is confirmed before
    FG voting, and the responsive honest votes are processed before the round
    closes; or
  \item a valid timeout certificate has already advanced~\(H\), and all
    responsive honest validators use its resulting height-\(H+1\) state.
\end{enumerate}
Suppose Assumption~\ref{ass:recovery-delivery} continues to hold and every later
step considered here satisfies Assumption~\ref{ass:goldfish}.  An honest
proposer at~\(H\) gives case~(a) unless it observes that case~(b) has already
occurred.  Then the common height eventually advances
\[
  H\longrightarrow H+1\longrightarrow H+2.
\]
After the SG votes for the second advance are delivered, there is a common
safe-confirmed prefix~\(C\).  Normally \(C=T_{H+1}\).  If a certificate has
already advanced~\(H+1\), let~\(C\) be the common target in the first later
round in which the target is confirmed before FG voting.  No honest validator
then adopts a conflicting anchor in any later round satisfying the same
assumptions.  Neither FG, SG, nor Goldfish can move the canonical
available chain to a branch that conflicts with~\(C\).
\end{theorem}

\begin{proof}
In case~(b), the first transition is already complete.  Since justification is
disabled at the nonjustifiable height~\(H\), it occurred by timeout.  In
case~(a), Lemma~\ref{lem:aligned-height-advance} gives a timeout quorum and the
round advances to~\(H+1\) by timeout, unless a certificate has already made
the same transition.  In either case, the next step starts from a common state
at height~\(H+1\).

If the next round meets the timing conditions stated after
Definition~\ref{def:round-close}, the first block at height~\(H+1\) is
given Goldfish available votes and confirmed before the next SG and FG votes.  All responsive
honest validators then have one chain, height, and current-height target, and
the delivered SG pre-votes make the target grade-0 clear before the FG action.
Otherwise wait for the next honest proposer after the transition state is
common.  Its proposal descends from that state by
Assumption~\ref{ass:recovery-delivery}.  Unless a certificate
has already advanced the height, its common height is~\(H+1\), and
Lemma~\ref{lem:honest-proposer-confirmation} confirms its common target before FG
voting.

If no intervening certificate has yet advanced~\(H+1\),
Lemmas~\ref{lem:aligned-height-advance}
and~\ref{lem:common-target-confirmation} advance it to~\(H+2\) and give
\(C=T_{H+1}\).  If a certificate has already made that transition, use the
next round meeting the conditions after Definition~\ref{def:round-close}, or
wait for the next honest proposer.  Its common target is confirmed before FG
voting by the round's timing or by Lemma~\ref{lem:honest-proposer-confirmation}; take
that target as~\(C\).  The two height advances are already complete in this
case, and the SG votes establish the same property for~\(C\).

By Lemma~\ref{lem:nonjustifiable-fg}, while
\(h_{\max}=H+1\), no \(J(H)\) exists and the cascade root is~\(F\), so FG
cannot interfere with the available walk.  Once \(h_{\max}=H+2\), any
eligible \(J(H+1)\) root is compatible with honest safe confirmations by
Lemma~\ref{lem:cascade-compatible}.  Lemma~\ref{lem:common-target-confirmation} makes
\(C\) common and safe-confirmed.  The Goldfish persistence assumption and the
SG pre-votes keep votes on descendants of~\(C\).
Lemma~\ref{lem:no-conflicting-anchor} then shows that fallback anchors cannot
choose a branch conflicting with~\(C\), while a conflicting pointed root would
require at least~\(\iota\) equivocation.  Hence the confirmation
established after the second advance persists for every later round in which
the assumptions hold.
\end{proof}

\begin{remark}[Number of honest proposers]\label{rem:honest-proposers}
One honest-proposer round is sufficient for one height advance, not necessarily
for a justification (Corollary~\ref{cor:honest-round-advance}).  Therefore two
honest-proposer rounds, the second beginning after the first transition state
is common, suffice for \(H\to H+1\to H+2\) if timely processing does not
continue from one round to the next.  They need not be consecutive wall-clock
rounds, but this gives no bound on the intervening rounds.

After the available chain is stable, an honest proposer at an ordinary height
with no previous honest votes produces a JC\@.  The following finality-piggyback
round does \emph{not} require an honest
first-slot proposer: it requires only that the common JC state be known before
the FG action and that the finality votes be included before the finality check
at round close.  Starting from a walk at~\(d-1\), or first entering a
nonjustifiable height, adds the rounds counted in
Theorem~\ref{thm:finality-liveness}.  Thus two consecutive honest proposers
followed by one arbitrary proposer suffice only when the second honest-proposer
round is already at an ordinary height with no previous honest votes.
\end{remark}

\begin{definition}[Maximum previous honest height]\label{def:previous-honest-votes}
Consider the first shared store with \(\Sigma.h_{\max}=H+2\) after the two
advances.  Let~\(m\) be an upper bound on:
\begin{enumerate}
  \item every state-height previously reached by an honest validator's stored
    walk; and
  \item every height at which an honest validator previously signed a target,
    timeout, or finality commitment.
\end{enumerate}
The bound covers everything before that shared store, including votes on
sibling branches before or after GST\@.  The condition \(H+2>m\) therefore
says that no honest validator had reached or signed at height~\(H+2\) before
the store became shared.  Honest walks in the shared store may still be at
either \(H+1\) or \(H+2\), as permitted by the height filter.  The next honest
proposer makes the choice common.
\end{definition}

\begin{lemma}[No pre-existing justification at \(H+2\) or above]\label{lem:no-preexisting-high-j}
Let~\(C\) be the persistent safe-confirmed prefix from
Theorem~\ref{thm:available-healing}.  Suppose the first shared store with
\(\Sigma.h_{\max}=H+2\) satisfies \(H+2>m\) in
Definition~\ref{def:previous-honest-votes}.  No pre-existing justification can redirect
the walk from~\(C\)'s chain:
\begin{enumerate}[label=(\roman*)]
  \item Definition~\ref{def:cascade-root} cannot select a \(J\) of height at
    most~\(H-1\);
  \item \(J(H)\) is impossible;
  \item every \(J(H+1)\) formed before or after the store became shared, and
    selected as a cascade root during the interval, is compatible with~\(C\); and
  \item no \(J(H+2)\) or higher can exist before this shared store.
\end{enumerate}
\end{lemma}

\begin{proof}
When \(\Sigma.h_{\max}=H+2\), Definition~\ref{def:cascade-root} can select a
justification only at height~\(H+1\), proving~(i).  Since~\(H\) is
nonjustifiable, \(J(H)\) is
impossible, proving~(ii).  A \(J(H+1)\) quorum contains an honest target vote.
If the justification is selected as a cascade root here, Assumption~\ref{ass:goldfish}
makes the confirmation that preceded that vote
compatible with~\(C\).
Target uniqueness excludes a conflicting certified target.  Any remaining
compatible height-\(H+1\) target lies on the same chain.  By
Lemma~\ref{lem:chain-target-uniqueness}, it is the same \(T_{H+1}\); hence every
\(J(H+1)\) names the same compatible target, proving~(iii).

Every quorum at height~\(x\) contains an honest vote, and an honest validator
can sign at~\(x\) only after a stored walk reaches~\(x\).  Since no honest
validator had reached or signed at any \(x\geq H+2\), no valid quorum
justification at such a height can predate the shared store.  This proves~(iv).
\end{proof}

\begin{definition}[Height bound after recovery]\label{def:height-bound-recovery}
After the available chain has stabilized, consider an honest-proposer round
whose proposal includes all blocks and still-valid honest votes received
before the proposal deadline.  Let~\(\Sigma\) be the shared store after processing
that proposal, its signed messages, and any preceding round state transition,
but before the available-vote action.  Write
\[
 d=\Sigma.h_{\max}.
\]
Assume also that \(d\geq x\) for every height~\(x\) that an honest validator
reached or signed at before this store.  No height transition occurs between
this store
and the round's available and SG walks.  The height filter therefore puts the
common decision height in \(\{d-1,d\}\), and no honest validator has a target,
timeout, finality vote, or lock above~\(d\) from before this store.
\end{definition}

\begin{lemma}[An ordinary height with no prior votes is justified]\label{lem:unused-ordinary-justify}
Let~\(u>d\) be a common ordinary height reached before any honest validator
has cast a target or timeout vote at~\(u\).  If its common target is confirmed
before FG voting and responsive honest votes of weight at least~\(q\) are
included before the round closes, then the current-height target~\(T_u\) is
justified and the chain advances to~\(u+1\).
\end{lemma}

\begin{proof}
No honest validator has an earlier height-\(u\) vote.  Every responsive honest
validator therefore votes for the same target~\(T_u\).  Their included weight
reaches~\(q\), so the state transition sets
\(\mathsf{JC}(u,T_u)\) before considering the simultaneously enabled timeout.
\end{proof}

\begin{lemma}[The next FG round finalizes the latest justification]\label{lem:later-finalization}
Suppose responsive honest validators of total weight at least~\(q\) signed the
target in \(\mathsf{JC}(u,T_u)\), and none of them timed out at~\(u\).  Before
their next FG action, suppose they have safe-confirmed a block whose state has
\((u,T_u)\) as the latest justified pair.  Suppose \((u,T_u)\) remains the
latest justified pair through collection of that action's finality fields; in
particular, no intervening state transition or branch change sets a newer
pair.  Every one of those validators then carries the finality commitment
\((u,T_u)\).  If those messages are included and processed at round close,
\(T_u\) finalizes there.  The round's first-slot proposer need not
be honest.
\end{lemma}

\begin{proof}
Each quorum signer previously targeted~\(T_u\), never timed out there, and now
knows the justification.  Since it is still the latest justified pair at the
action, the signer satisfies Definition~\ref{def:finality-vote} and records its
lock on~\(T_u\), regardless of whether its new current-height decision is
target, timeout, or empty.  These commitments add signer weight at least~\(q\)
to~\(P\).  Since finality is checked first at round close, it sets
\((h_F,F)=(u,T_u)\) before a simultaneous target quorum can set a newer
justification and reset~\(P\).  This is why the lemma requires \((u,T_u)\) to
remain the latest justification until all finality fields are collected.
\end{proof}

\begin{theorem}[Finality after recovery]\label{thm:finality-liveness}
Assume \(H+2>m\) as in Definition~\ref{def:previous-honest-votes} and
Assumption~\ref{ass:recovery-delivery}.  After the available chain has stabilized,
let~\(d\) be as in Definition~\ref{def:height-bound-recovery}.  Consider five
consecutive rounds \(r_1,\ldots,r_5\).  In each round, responsive honest
validators of total weight at least~\(q\) complete the available, SG, and FG
steps, and their votes are included by that round's close.  Whenever a round
advances the height, before the next round's SG and FG votes every responsive
honest validator has the same resulting state and has safe-confirmed a block
carrying that state.  In particular, if the advance forms a justification, it
is the latest justification in that state.  The first
round has an honest first-slot proposer and a common proposal before the
available-vote action.  Later rounds need not have honest proposers.

Counting \(r_1\) as the first round, finalization occurs by \(r_5\).  It occurs
by \(r_4\) if height~\(d+1\) is ordinary.  It occurs by \(r_3\) only when the
first common walk has height \(k=d\) and height~\(d+1\) is ordinary.

Without five consecutive rounds satisfying these conditions, finality still
resumes if an honest proposer later reaches a common ordinary height~\(u>d\)
before any honest validator votes at~\(u\), its target is confirmed before FG
voting, and the target votes are included at round close.  Before the next FG
action, every responsive honest validator must safe-confirm a block carrying
the resulting justification.  The justification must remain latest until that
action's finality fields are collected, and the included finality votes must
have total weight at least~\(q\).  This
eventual result has no fixed round bound.
\end{theorem}

\begin{proof}
Lemma~\ref{lem:no-preexisting-high-j} excludes every pre-existing justification that could
move the walk to a conflicting chain, and Lemma~\ref{lem:no-conflicting-anchor}
keeps later walks on the stabilized chain.  In the shared store, \(d\) is the
value of~\(h_{\max}\) and bounds all heights previously reached or
signed at by an honest validator.  The height filter and the honest proposal
therefore give a common first decision height \(k\in\{d-1,d\}\).
Lemmas~\ref{lem:honest-proposer-confirmation}
and~\ref{lem:aligned-height-advance} advance~\(k\).  If \(k=d-1\), the same
argument in the next counted round advances~\(d\), so the common chain enters
height~\(d+1\) after at most two rounds.

No honest vote at height~\(d+1\) predates the shared store.  The assumptions on
the counted rounds give validators a common confirmed target before the first
FG vote at that height.
If the height is ordinary, Lemma~\ref{lem:unused-ordinary-justify} justifies its target.
If it is nonjustifiable, the timeout rule advances it.  Consecutive heights
cannot both be multiples of~\(K\geq2\), so the next height is
ordinary and Lemma~\ref{lem:unused-ordinary-justify} applies there.  The new
justification is carried by a common safe-confirmed block before the next FG
action.  Every
justification quorum contains honest weight, and~\(d\)
bounds all heights previously reached or signed at by honest validators.
Therefore no branch state with a higher justification can predate the shared
store.
Because the height rules run only at round close, they set no newer
justification before the finality fields are collected, so
Lemma~\ref{lem:later-finalization} finalizes the latest justified target.

If \(k=d\), entry into~\(d+1\), its ordinary justification, and its finality
use rounds 1, 2, and 3; an intervening nonjustifiable~\(d+1\) adds one round.
If \(k=d-1\), the extra round gives the four- and five-round cases.  For a
later ordinary height~\(u>d\), the theorem's stated conditions satisfy
Lemmas~\ref{lem:unused-ordinary-justify}
and~\ref{lem:later-finalization}, which give eventual finality without a numerical
bound.
\end{proof}

\begin{theorem}[Healing]\label{thm:eventual-healing}
Under Assumptions~\ref{ass:height-weights},
\ref{ass:recovery-weights}, \ref{ass:goldfish},
\ref{ass:honest-proposer}, and~\ref{ass:recovery-delivery}, suppose recovery reaches
a height~\(H\) that is nonjustifiable on every valid branch, including a
branch revealed later.  Suppose a shared store has
\(\Sigma.h_{\max}=H\), and a future round has an honest first-slot proposer
using that store.  Its common height is in \(\{H-1,H\}\).  If it is~\(H-1\),
suppose either the next round satisfies the timing conditions after
Definition~\ref{def:round-close}, or a later honest proposer uses the resulting
shared state.  Corollary~\ref{cor:height-catchup} then brings the shared chain
to~\(H\).  If the common height is already~\(H\), no preliminary round is
needed.  In either case Theorem~\ref{thm:available-healing} applies.  Suppose also that the
first shared store with \(h_{\max}=H+2\) satisfies \(H+2>m\) as in
Definition~\ref{def:previous-honest-votes}.  For finality, additionally suppose that
the conditions of Theorem~\ref{thm:finality-liveness} hold.  Then:
\begin{enumerate}
  \item the common height advances through \(H\to H+1\to H+2\);
  \item after the SG votes for the second advance are delivered, the available
    chain never switches to a chain conflicting with the resulting
    safe-confirmed prefix during any
    later action in which Assumption~\ref{ass:goldfish} holds; and
  \item the target at a later ordinary height with no previous honest votes is
    justified and finalized.
\end{enumerate}
The accountable-safety theorem does not use the synchrony or participation
assumptions.
\end{theorem}

\begin{proof}
Corollary~\ref{cor:height-catchup} gives the first case of
Theorem~\ref{thm:available-healing} immediately or after one preliminary round,
unless a certificate has already given its second case.  That theorem proves
the two advances and stability of the available chain.  Lemma~\ref{lem:no-preexisting-high-j}
excludes any remaining hidden justification that could affect fork choice.
Theorem~\ref{thm:finality-liveness} then gives a new justification and
finalizes it within its stated round bound.
\end{proof}

\begin{corollary}[Recurring healing]\label{cor:recurring-healing}
If every later finite disruption is followed by another uninterrupted window
satisfying the theorem's premises at a new height that is nonjustifiable on
every valid branch, then another block is finalized after each disruption.
Finitely many disruptions followed by one stable period are the special case
in which finality resumes.  Continued ordinary finality additionally uses the
usual delivery, participation, and inclusion assumptions, which are not proved
here.
\end{corollary}

\begin{proof}
Apply Theorem~\ref{thm:eventual-healing} after each finite disruption.  Each
newly finalized target has a strictly greater height than the preceding one,
so every application strictly advances~\(h_F\).
\end{proof}

\section{Incentives and empty decisions}\label{sec:incentives}

The safety and healing theorems do not depend on rewards or the inactivity
leak.  The only point relevant to rewards is that an empty decision
sets no timeout bit and cannot exempt its signer from a stall penalty.
Otherwise all validators could vote empty, avoid the penalty, and prevent
height advance forever.

This paper does not prove a quantitative lower bound for the inactivity leak.
A deployment may pay
FG target and finality participation once per round or once per height.  The
design must specify eligible validators, rounding, fork-transition accounting,
and which validator balances and epoch are used to calculate rewards.  Paying
target/finality flags once per height avoids repeatedly rewarding the same
stalled checkpoint, but the practical gain is limited because the inactivity
leak quickly removes positive rewards during a stall.  The liveness proof is
independent of that choice.

\section{Relation to the executable specification}\label{sec:refinement}

The proof places SG and FG votes after Goldfish confirmation, while an
executable protocol may spread these duties across a round.  Such a schedule
must preserve the same order for every vote used in the proof.  The validator
set and weights used at a height must remain fixed across forks.  Checkpoint
sync must authenticate \texttt{current\_height\_target} together with the
current height and finalized state.  A valid vote that refers to an unknown
block must be queued and replayed when the block arrives.  The companion note
\texttt{paper-spec-status.md} records the remaining specification details.

\section{Conclusion}\label{sec:conclusion}

Fresh Simplex uses fixed weights at each height, E1 and E2, and one
target-participation bit and one timeout bit per validator.  A proposer may
point to a root supported in the previous round.  Goldfish first confirms the
proposal chain; SG and FG then vote from that confirmed chain.

The recovery argument begins at a height~\(H\) that is nonjustifiable on every
valid branch, including a branch revealed later.  Advancing through
\(H\to H+1\to H+2\) makes the confirmed available prefix stable.  Excluding an
old justification at height~\(H+2\) or above requires \(H+2>m\), where \(m\)
bounds every height previously reached or signed at by an honest validator.
Finality additionally requires a later ordinary height with no previous honest
votes, timely inclusion of its target votes, and a following round in which
that justification remains latest until the finality votes are included.

The result assumes responsive honest validators of total weight at least~\(q\),
fixed weights during the
recovery interval, data and ancestry availability, compatible and persistent
Goldfish confirmations, an honest proposer to make the chain and target
common, and timely delivery and inclusion for the rounds being counted.  With
only recurring honest proposers the height still advances, but there is no
fixed round bound and finality needs a later ordinary height above all previous
honest votes before any honest validator votes there.

\end{document}
